\section*{Associate Editor}

\comment{
The author has effectively addressed some of the comments from the initial submission.
Two new referees have reviewed the paper and give several suggestions for further revision.
In addition, the authors should consider the following points.
}

\reply{
	We thank the associate editor for the careful consideration of our manuscript. We have addressed their points and those of the two reviewers and improved the manuscript as a result. Our responses are presented in a point-by-point fashion below.
}

\comment{
The MFPCA approach was described more clearly in the revision, which leads to an
important point about visualization. In equation $(2)$, the eigenfunctions $\phi_k(s,t) \in \mathbb{R}^2$
are bivariate, with one component corresponding to made and the other to missed
shots. Doing MFPCA instead of marginal FPCA is nice since variability in made and
missed shot locations are taken together, yielding a single set of scores for each player
to cluster. However, there is no indication in the figures regarding how much each
eigenfunction loads onto made and missed shots. In other words, it is possible that
$\phi_1(s,t)$ loads mostly onto made shots and almost none onto missed shots, or vice versa,
but there is no way to tell this from the current analysis because the normalization (a
linear transformation, I assume) is applied separately to every panel in each figure. It
would be more informative to apply this transformation to the combined data across
made and missed shots - that is, normalization is applied to each column. Also, it
would be helpful to report how much each component contributes to the variability
explained in made/missed shots separately, in addition to the joint figures currently
reported.
}

\reply{
  We agree with this comment. The separate normalization of each panel is useful for visualizing spatial patterns, but hides the relative magnitude of the made-shot and missed-shot components. We have therefore change the visualization. The normalization is now applied jointly to the made-shot and missed-shot surfaces within each principal components. More precisely, for each column of the figures, we computed the maximum value of the densities and divided the made-shot and missed-shot densities by this value. We do not see much difference with the previous version, indicating that the load of each eigenfunction onto made and missed shots is balanced.


Following \cite{ramsayFunctionalDataAnalysis2005}, we reported $\|\phi_{k \text{missed}}\|^2$ and $\|\phi_{k, \text{made}}\|^2$ to show how much each component contributes to the variability explained in made/missed shot. These quantities show the relative contribution in each dimension. This is a relative measure because
$$ \|\phi_{k,\text{missed}}\|^2 + \|\phi_{k, \text{made}}\|^2 = 1,$$
by the orthonormality property.
}

\comment{
Although the author has argued that lack of consensus in a suitable spatial grid leads to subjectivity in previous analyses, there is no apparent justification for the choice of a $201 \times 201$, equally spaced grid for the MFPCA. In addition, the density estimation induces further subjectivity by choice of a bandwidth. Is there a way to assess sensitivity of the results and interpretations to these choices? It seems the choice is no less arbitrary for density representations compared to previous methods.
}

\reply{
  We have added a sensitivity analysis to the revised manuscript. The $201 \times 201$ grid is used to evaluate the smooth kernel density estimates and approximate the MFPCA eigenfunctions. It does not define shooting regions, impose boundaries or create discrete shooting categories. To assess the numerical effect of this choice, we repeated the complete estimation procedure using two coarser grids, $51 \times 51$ and $101 \times 101$, and one finer grid, $401 \times 401$.
  The dominant spatial structures were stable across these resolutions. The mean functions and first eigenfunctions were very similar to those obtained with the $201 \times 201$ grid, and the $401 \times 401$ results were nearly indistinguishable from the main analysis. Some differences appeared for the coarsest $51 \times 51$ grid and for lower-variance components, which is expected because later components are more sensitive to numerical approximation. These results support the use of the $201 \times 201$ grid as a sufficiently fine approximation while avoiding unnecessary computational cost.

  We also added a bandwidth sensitivity check based on cross-validation. For each player and shot outcome, shot locations were rescaled to $[0,1] \times [0,1]$, and Gaussian kernel density estimators were fitted over a $20$-point logarithmic grid of candidate bandwidths from $10^{-3}$ to $1$. Each candidate bandwidth was evaluated using five-fold cross-validation, with the selected bandwidth maximizing the average held-out log predictive density. The resulting densities showed spatial patterns similar to those obtained with Silverman's rule, indicating that the main features of the analysis are not driven by the particular rule-of-thumb bandwidth used in the primary results.

  We have therefore clarified that the grid and bandwidth are numerical implementation choices in the functional-density representation. Unlike fixed spatial bins or manually defined shot zones, they do not determine the substantive categories being analyzed, and the added sensitivity checks show that the main interpretations are stable to reasonable alternatives.
}

\comment{
Related to the previous point, it is clear that densities are able to examine patterns without the influence of volume. However, how would removing volume information prove valuable in making performance-related decisions in terms of coaching, scouting, or training, for example? Are there other solutions that could be more effective, such as stratifying by volume rather than removing it altogether? Please address this more completely in the discussion.
}

\reply{
{\color{purple}\textbf{Ed To Do}: Double check and either reference discussion or add part.}
  Our density-based representation is intended to characterize where a player tends to shoot, conditional on the player having taken shots, rather than how often the player shot shoots overall. This is useful in making performance-related decisions when the goal is to compare spatial shooting tendencies independently of playing time, number of games or usage. For example, two players with very different total shot counts can still have similar spatial tendencies. Our approach allows this to be detected.

  \textcolor{blue}{We should add something in the discussion.}

  We could add intensity surfaces as a third coordinate in the vector-valued stochastic process. \textcolor{blue}{I am not sure it is a good idea to stratify by volume, because the estimation of the densities depend on the number of shots we have for the player. So, players with less shots will have their densities less accurately estimated. Also, how do we choose the strate?} Alternatively, shot volume could also be included as an additional feature in the clustering step.
}

\comment{
The choice of $k= 5$ is clearly not data-driven. What would a true data-driven clustering look like? Certainly there are more than 5 shot pattern styles, so what information can we really gather from the clustering beyond observing that they clusters with $k= 5$ are poorly aligned with player position? Also, although several metrics on the clustering are reported (ARI, Silhouette score), the relevance of the values is not discussed or their bearing on the reliability and impact of the clustering results.}

\reply{
  We agree that the choice of $k = 5$ is not data-driven. Its purpose is to provide a controlled comparison with the five conventional NBA position groups, not to claim that there are exactly five natural shooting styles. We have clarified this point in the revised manuscript and added a new \emph{Data-driven clustering} subsection.

  In that subsection, a fully data-driven version of the clustering keeps the same MFPCA score representation, the same $k$-medoids algorithm and the same distance matrices, but selects the number of clusters by maximizing the average Silhouette coefficient over $k = 2, \ldots, 10$. With equal weights for the four MFPCA scores, the maximum Silhouette coefficient is $0.18$ and is obtained for $k = 3$. The three clusters contain $41$, $57$ and $75$ players, with Norman Powell, Coby White and Dennis Schröder as medoids. With variance-based weights, the maximum Silhouette coefficient is $0.54$ and is obtained for $k = 2$. The two clusters contain $99$ and $74$ players, with Julius Randle and Jordan Poole as medoids. Because the first component explains most of the variation, the variance-weighted data-driven clustering is mainly a split along the first score.

  The negative Silhouette coefficients for the NBA position labels ($-0.05$ and $-0.1$, depending on the distance matrix) reinforce that conventional positions are not compact or well-separated groups when evaluated using shooting-pattern distances. For the five-cluster $k$-medoids solutions, the equal-weighted Silhouette coefficient of $0.04$ indicates weak separation and supports interpreting those clusters as descriptive summaries rather than as reliable evidence of five isolated shooting types. The variance-weighted Silhouette coefficient of $0.4$ indicates a clearer partition, but one driven mainly by the dominant components. Thus, the $k = 5$ analysis is useful for comparison with NBA positions, while the data-driven analysis suggests that the strongest intrinsic grouping in the MFPCA score space is coarser.
}
